\begin{textsample}{POS Dim 1 – ro – Score 41.00 – ro/ro\_em\_15.txt}  \label{ex:f1_pos_147}
1.6.2.
Trabalho pedagógico com as metodologias ativas no ensino médio A expansão da era digital e a demanda por recursos tecnológicos presentes na sociedade da informação, do conhecimento ou da aprendizagem propõem \textbf{inúmeros} desafios ao trabalho pedagógico na \textbf{contemporaneidade} ( ALVES E SOUSA, 2016 ).
Dentre os desafios, a nova cultura de aprendizagem ativa para atender o perfil de estudante \textbf{revela} a importância da preparação do professor para o uso e a atuação com as ferramentas tecnológicas em diferentes espaços de aprendizagem e nas diferentes áreas de conhecimento contempladas no currículo para Educação Básica.
Outro ponto importante \textbf{diz} respeito \textbf{que} ao cenário educacional, devido à falta de interesse do estudante, o distanciamento da sala de aula e os \textbf{métodos} tradicionais de ensino \textbf{que} \textbf{desencadeiam} importantes reflexões acerca da significação da aprendizagem e sobre a inovação dos processos educativos.
Concomitante, entende -se \textbf{que} é \textbf{crucial} a formação continuada dos professores pautada em pressupostos epistemológicos \textbf{inerentes} a \textbf{uma} prática educativa mediadora, crítica e reflexiva na qual o estudante é o centro da aprendizagem ( FIALHO E MACHADO, 2017 ), o \textbf{que} corrobora com a relevância e \textbf{aplicabilidade} das metodologias ativas no Ensino Médio.
Nesse contexto, o trabalho pedagógico pautado em metodologias ativas no Ensino Médio é reflexo de \textbf{um} processo de metamorfose existente na sociedade \textbf{contemporânea} sob o qual tem sido vislumbrada a necessidade emergente de repensar o ensino, os \textbf{métodos} e as estratégias utilizadas pelo professor em sua prática educativa.
Como pressupostos iniciais, constituem fatores \textbf{que} culminam para emergência na transformação do ensino na Educação Básica, as mudanças globais no estilo de vida, no modo de \textbf{viver}, conhecer, ser, pensar, agir e aprender do estudante, e consequentemente, no ato de \textbf{educar} e formar essa nova geração.
O estudante da \textbf{contemporaneidade} requer \textbf{uma} formação \textbf{que} contemple as exigências da era digital, da comunicação e do conhecimento, sendo assim, mostra -se como essencial propiciar experiências educativas \textbf{que} o envolva como \textbf{pesquisador} e investigador para resolver problemas concretos \textbf{que} ocorrem no cotidiano de suas vidas, na sua comunidade e na sociedade de forma geral.
A aprendizagem precisa ser significativa, desafiadora, problematizadora e instigante, a ponto de mobilizar o estudante e o grupo a buscar soluções possíveis para serem discutidas e concretizadas à luz de referenciais \textbf{teóricos} práticos ( MORAN, MASETTO E BEHRENS, 2013 p.83 ).
Em vista disso, o referencial curricular do Estado de Rondônia aponta o campo das metodologias ativas com sendo \textbf{um} caminho \textbf{metodológico} e estratégico relevante e fundamental para reflexão na ação e sobre a ação ( SCHON, 2000 ) no \textbf{que} \textbf{diz} respeito a re(criar) experiências educativas significativas para o estudante do Ensino Médio nas áreas, entre as áreas de conhecimentos e na parte flexível do currículo. \textbf{Ademais}, defende a importância da construção de \textbf{uma} prática educativa fundamentada na Educação por competência \textbf{que} forneça bases para \textbf{que} o professor possa compreender as mudanças nas concepções do aprender, planejar, ensinar e comunicar na sociedade da informação, do conhecimento e da aprendizagem.
Atrelada a isso, a produção do saber nas áreas do conhecimento demanda pelo desenvolvimento de ações \textbf{que} conduzem o professor e o estudante na busca dos processos de investigação e pesquisa com vistas a abrir caminhos coletivos para produção de conhecimento de forma integrada por parte do aluno e professor.
Nesse âmbito “ é necessário \textbf{que} o estudante ultrapasse o papel passivo de escutar, ler, decorar e de ser repetidor fiel dos ensinamentos do professor, para se tornar criativo, crítico, \textbf{pesquisador} e atuante para produzir conhecimento. ” ( MORAN, MASETTO E BEHRENS, 2013, p. 77 ).
Cabe destacar ainda \textbf{que}, as novas ferramentas tecnológicas presentes no mundo \textbf{contemporâneo}, além de \textbf{exigir} a ressignificação da prática educativa em sala de aula e o uso de novos ambientes de aprendizagem, também perpassam pela compreensão da aprendizagem situada em bases \textbf{teóricas} e práticas.
Desse modo, é relevante por parte do professor, a adoção de \textbf{uma} concepção crítica e reflexiva da ação docente \textbf{que} considera as especificidades dos estudantes no \textbf{que} \textbf{diz} respeito às mudanças sociais, políticas, financeiras, tecnológicas, e a peculiaridade regional e local.
Referente à perspectiva de aprendizagem significativa, a utilização de metodologias ativas no Ensino Médio favorece a autonomia, a curiosidade, a tomada de decisão individual e coletiva do estudante, e ainda, o protagonismo para o desenvolvimento de atividades \textbf{que} \textbf{permeiam} práticas sociais diferenciadas em conformidade com os contextos em \textbf{que} o estudante está inserido ( MORAN, MASETTO E BEHRENS, 2013 ; FIALHO E MACHADO, 2017 ).
Além disso, do ponto de vista da formação integral do estudante, conforme as competências gerais da BNCC, o trabalho pedagógico com as metodologias ativas contribui de forma direta no desenvolvimento social, conhecimento, criatividade, capacidade crítica reflexiva, autoconfiança, comunicação, cooperação, trabalho de equipe e empatia do estudante.
Sob essa perspectiva, as situações e os espaços de aprendizagem devem ser impulsionados pela curiosidade, pelo interesse, pela \textbf{crise}, pela \textbf{problematização} e pela busca de soluções possíveis para aquele momento histórico com a visão de \textbf{que} não há respostas únicas, absolutas e inquestionáveis ( MORAN, MASETTO E BEHRENS, 2013 ).
Nessa \textbf{vertente}, é necessário, segundo Freire ( 1996 ), \textbf{que} o docente perceba o \textbf{educando} \textbf{enquanto} sujeito \textbf{sócio-histórico} e cultural do ato de conhecer e amplie suas possibilidades educativas respeitando a coerência entre o saber-fazer e o saber ser.
Em suma, é relevante a construção de ambientes favoráveis à produção de conhecimento e desenvolvimento da autonomia do estudante.
Sendo assim, experiências como atividades realizadas em grupos com mais de \textbf{um} professor de classe acompanhando a execução de \textbf{tarefas}, realização de projetos, soluções de problemas reais e estudo de caso são estratégias \textbf{que}, se bem conduzidas, podem gerar \textbf{uma} verdadeira inovação pedagógica.
Cabe evidenciar ainda, \textbf{que} o ato de inovar em espaços de aprendizagem e na própria sala de aula amplia as possibilidades para estabelecer relações significativas entre os diferentes saberes, de maneira progressiva e mais elaborada, bem como converte a escola em lugar mais democrático, atrativo e estimulante ( DARO, 2018 ).
Em suma, é essencial no trabalho pedagógico, a mobilização e articulação de saberes considerando as experiências reais e/ou simuladas do estudante ( SOUZA, VILAÇA E TEIXEIRA, 2020 ) no intuito de oportunizar as condições para \textbf{que} possa solucionar problemas e demandas complexas presentes no seu cotidiano.
Sendo assim, no campo da \textbf{aplicabilidade} e do trabalho com as metodologias ativas nos espaços de aprendizagem, considerando a diversidade de estratégias, Camargo ( 2018b ) propõe, sob a perspectiva da sala de aula inovadora, estratégias para \textbf{fomentar} o aprendizado ativo, dentre as quais se destacam : a aprendizagem \textbf{baseada} em problemas, em projetos, a aprendizagem entre pares ou times, design thinking, gamificação, dentre outros.
Em relação ao trabalho com Design Thinking, Cavalcanti e Filatro ( 2016 ) afirmam \textbf{que} \textbf{enquanto} estratégias de aprendizagem correspondem a “ \textbf{um} modo de pensar ” \textbf{métodos} e estratégias para produção criativa de soluções inovadoras e promoção da \textbf{centralidade} do estudante no processo educativo.
Para tanto, constituem elementos chaves para as necessidades e as experiências do estudante, desde \textbf{que} imbuídos em três momentos : a inspiração, a ideação e a implementação.
Nesse contexto, a figura 5 \textbf{ilustra} \textbf{um} organograma com as etapas para criação de \textbf{uma} atividade com base na estratégia do Design Thinking, seguindo a técnica HCD Toolkit \textbf{que} \textbf{diz} respeito a hear ( ouvir ) ; create ( criar ) ; deliver ( entregar ).
Mediante à figura 4, cabe ressaltar \textbf{que} o Design Thinking na condição de estratégia propicia ao estudante criar opções e fazer escolhas \textbf{que} contribuam na produção de novas ideias e na projeção de soluções inovadoras.
Outro ponto de destaque, \textbf{diz} respeito \textbf{que} por meio da prototipagem, é possível tornar as ideias tangíveis para \textbf{que} sejam testadas e \textbf{que} permitam a discussão e/ou análise dos pontos \textbf{fortes} e fraquezas no desafio proposto.
Outra aplicação prática das metodologias ativas proposta por Flora ( 2015 ) é a gamificação.
Para a autora, a gamificação na Educação é \textbf{uma} estratégia de games ( jogos ), utilizada para promover o engajamento no ambiente de aprendizagem \textbf{que} por sua funcionalidade, \textbf{revela} -se como sendo \textbf{um} recurso \textbf{que} não pode faltar na caixa de ferramentas do professor, \textbf{uma} vez \textbf{que} auxilia no alcance dos objetivos de aprendizagem de forma divertida e propõe como foco \textbf{central} o desenvolvimento de estratégias para resolução de problemas \textbf{que} contribuem efetivamente para a aprendizagem significativa do estudante.
Nesta \textbf{vertente}, para o planejamento do trabalho pedagógico com estratégias gamificadas o organograma abaixo ( figura 6 ) orienta com base em Flora ( 2015 ) aspectos relevantes \textbf{que} devem ser observados quando a meta ou propósito do game ( jogo ) é a aprendizagem.
Figura 6.
Organograma explicativo dos aspectos relevantes para o planejamento de \textbf{uma} atividade gamificada para o estudante do Ensino Médio no Estado de Rondônia. | Aspecto | Descrição | |---|---| | Meta | Motiva a ação do estudante ; a execução de \textbf{tarefas} ; a solução de problemas ; as mudanças de comportamento ; e o grau de dificuldade. | | Engajamento | Compartilhar o conhecimento ; contribuir para alcançar resultado ; criar algo \textbf{que} seja \textbf{interessante} e envolvente. | | Resolução de problemas | Oportuniza pensar sobre \textbf{um} problema ou atividade do dia a dia e convertê -lo em \textbf{uma} atividade \textbf{que} contenha os elementos do jogo ( competição, cooperação, exploração e premiação ). | | Pensamento de jogo | \textbf{Desperta} o interesse e prazer em jogar ; e a \textbf{percepção} de \textbf{que} participa e interage na construção de algo \textbf{enquanto} interagem. | | Diversão | Intensifica por meio de games a natureza competitiva, competitiva, o jogo com pessoas diferentes, se aceita o feedback. | Fonte : Adaptado de Flora ( 2015 ) Perante o organograma apresentado na figura 6, percebe -se a necessidade de articulação dos elementos do game, o pensamento de jogos e técnicas de design de forma dinâmica para motivar a ação, interação e o feedback no game realizado.
Sob a perspectiva da inovação e progressão da aprendizagem, destaca -se a participação e o envolvimento do estudante em atividades gamificadas com o uso de QR Code, pois oportunizam ao mesmo decifrar pistas e missões contidas nesses códigos, ao mesmo tempo em \textbf{que} auxilia no desenvolvimento da criatividade, autonomia, diálogo, na resolução de problemas e pode ser usado em componentes curriculares distintos. \textbf{Face} ao \textbf{exposto}, chama a atenção para o fato de \textbf{que} a estratégia de gamificação no campo da Educação não é aplicada apenas com o uso de tecnologias, pois existem formas mais primitivas e comuns ao ambiente escolar \textbf{que} o professor sequer imagina \textbf{que} possa ser atividade gamificada.
Tal afirmação encontra fundamentação quando se analisam os jogos \textbf{populares}, jogos de tabuleiros e jogos pedagógicos realizados em sala de aula para diversão, engajamento e alcance de metas.
Como exemplo disso, no planejamento de \textbf{uma} atividade gamificada recomenda -se : a criação de \textbf{um} roteiro a ser cumprido, juntamente com a as missões a serem \textbf{desvendadas} pelos estudantes e, após a experiência gamificada, o professor pode realizar discussões com apresentação de ideias no quadro ou esquemas em papéis, incentivando os estudantes a criar de maneira colaborativa.
Por fim, no intuito de ampliar as possibilidades educativas para o trabalho pedagógico do professor com os estudantes do Ensino do Estado de Rondônia optou -se por apresentar no quadro 2 \textbf{uma} proposição de \textbf{detalhada} com exemplos de estratégias de metodologias ativas para o trabalho pedagógico e a mobilização dos saberes dos estudantes de forma ativa e significativa.

% matched lemmas: ademais, aplicabilidade, basear, central, centralidade, contemporaneidade, contemporâneo, crise, crucial, desencadear, despertar, desvendar, detalhar, dizer, educar, enquanto, exigir, exposto, face, fomentar, forte, ilustrar, inerente, interessante, inúmero, metodológico, método, percepção, permear, pesquisador, popular, problematização, que, revelar, sócio-histórico, tarefa, teórico, um, vertente, viver
\end{textsample}
